\abstract
\keywords{ОБОГАЩЕНИЕ ТОВАРНЫХ ДАННЫХ, КАСКАДНАЯ АРХИТЕКТУРА, СЛАБЫЙ НАДЗОР, OPEN FOOD FACTS, БАЙЕСОВСКАЯ СЕТЬ, БОЛЬШАЯ ЯЗЫКОВАЯ МОДЕЛЬ, КАЛИБРОВКА, ИЗВЛЕЧЕНИЕ ПРИЗНАКОВ, ДОМИНИРОВАНИЕ ПО ПАРЕТО, АУДИТ-ЦИКЛ, ВЕКТОРНЫЕ ПРЕДСТАВЛЕНИЯ}

\textbf{Объём работы.} 91 страница основной части, 4 главы, 9 рисунков, 22 таблицы, 52 источника (16 российских и 36 зарубежных).

\textbf{Объект исследования.} Процесс автоматизированного обогащения товарных каталогов интернет-магазинов с использованием открытых источников данных.

\textbf{Предмет исследования.} Гибридная каскадная архитектура с ансамблевой оценкой уверенности классификаторов и адаптивным обращением к большой языковой модели как запасному слою.

\textbf{Цель работы.} Разработать гибридную каскадную систему автоматизированного обогащения товарных данных для интернет-магазина, обеспечивающую высокую точность извлечения категориальных атрибутов при сокращённой стоимости обращения к большой языковой модели и поддерживающую развёртывание на открытых моделях.

\textbf{Методы.} Многоязычные векторные представления текста на базе модели paraphrase-multilingual-MiniLM-L12-v2 (SBERT, поддержка 50+ языков), классификатор XGBoost с регуляризацией и поатрибутными порогами уверенности, байесовская сеть со структурой, найденной алгоритмом Hill Climb по критерию BIC (библиотека pgmpy), доверительные интервалы с кластеризацией по брендам через ресэмплированный бутстрэп, бакетирование числовых нутри-полей, согласие трёх независимых больших языковых моделей (Sonnet 4.5 + GPT-4o + Gemini 2.5 Flash) для построения эталонной разметки на текстово-семантических атрибутах, слепая верификация моделью Claude Opus 4, разбиение на обучающую и тестовую части без пересечения товарных кодов (code-disjoint), проверка гипотезы, зафиксированной до получения результатов с поправкой Бонферрони на множественные сравнения.

\textbf{Результаты.} Реализована и эмпирически проверена гибридная каскадная архитектура из пяти слоёв (категорайзер Слой 0 — регулярные выражения — классификатор XGBoost на векторных представлениях SBERT — байесовский валидатор плотностей — запасной слой большой языковой модели) на трёх продуктовых категориях открытого каталога Open Food Facts (паста, шоколад, сыры). На тестовой выборке без пересечения товарных кодов (code-disjoint, новые SKU) размером 3257 ячеек по 20 парам «категория-атрибут» итоговая конфигурация достигает 91,1 \% сквозной точности при обращении к большой языковой модели на 7,1 \% ячеек в режиме «каскад + gemini-flash»; точность только-каскадной части (без учёта непокрытых ячеек) составляет 94,8 \% (macro-F1 0,899). На консервативной нижней оценке с независимой экспертной разметкой моделью Claude Opus 4 (n=615) — 87,5 \% сквозной точности. При развёртывании на собственных мощностях с открытой моделью gpt-oss-120b — 90,6 \% точности, что на 21,3 п. п. выше прямого использования gpt-oss-120b (69,3 \%). Каскад без запасного слоя покрывает 96,0 \% ячеек. Достигнуто доминирование по Парето разработанной системы над прямым использованием большой языковой модели по двум осям одновременно — точности и стоимости (+7,3 п. п. точности и сокращение стоимости в 333 раза при той же экономии по сравнению с прямым использованием Claude Sonnet 4.5; архитектурное сокращение стоимости в 14 раз обеспечивается переводом 92,9 \% ячеек на дешёвые слои каскада). Гипотеза, зафиксированная до получения результатов, о превосходстве обучаемого стоимостно-чувствительного маршрутизатора над статическим правилом отклонена после поправки Бонферрони на трёх бюджетах. Реализован замкнутый цикл «аудит — правки экстрактора — переобучение — измерение» с кросс-категорийной репликацией 3 из 3. Разработана демонстрационная система из трёх компонентов (шлюз API на Go, ML-сервис на Python/FastAPI, веб-интерфейс), реализующая стабильный контракт REST API.

\textbf{Использование результатов.} Разработанная архитектура применима в системах управления товарной информацией (PIM-системах) интернет-магазинов. Внедрение системы позволяет: повысить полноту автоматического заполнения карточек товаров до 91,1 \% (сквозная точность с учётом маршрутизации; 94,8 \% по каскадной части как верхняя граница при идеальной маршрутизации категории), сократить стоимость обработки в 333 раза по сравнению с прямым использованием Claude Sonnet 4.5 (комбинированное сокращение, включающее архитектурный вклад каскада в 14 раз и удешевление модели запасного слоя), сократить ручную нагрузку оператора каталога в 10–14 раз за счёт перевода в режим выборочной верификации, обеспечить развёртывание на собственных мощностях без отправки партнёрских данных во внешние API больших языковых моделей.

